\lysh{13174}{4}

\begin{alist}
\item Η παράσταση Α ορίζεται όταν:
\[|x| - 1 \ne 0 \Leftrightarrow |x| \ne 1 \Leftrightarrow x \ne \pm 1\]
και η παράσταση Β όταν:
\[|x| - 2 \ne 0 \Leftrightarrow |x| \ne 2 \Leftrightarrow x \ne \pm 2.\]

\item Έχουμε: $-x^2 + 4|x| - 3 = -|x|^2 + 4|x| - 3$. Θέτουμε $\omega = |x|$, οπότε
\[-|x|^2 + 4|x| - 3 = -\omega^2 + 4\omega - 3,\]
που είναι τριώνυμο με ρίζες των οποίων το άθροισμα είναι $S = -\dfrac{4}{-1} = 4$ και το γινόμενο
είναι $P = \dfrac{-3}{-1} = 3$, οπότε $\omega_1 = 3$ και $\omega_2 = 1$. Άρα
\[-\omega^2 + 4\omega - 3 = -(\omega - 1)(\omega - 3)\]
και
$-x^2 + 4|x| - 3 = -(|x| - 1)(|x| - 3)$.

Συνεπώς: $\text{A} = \dfrac{-x^2 + 4|x| - 3}{|x| - 1} = \dfrac{-(|x| - 1)(|x| - 3)}{|x| - 1} = 3 - |x|$.

Για την παράσταση Β έχουμε:
\[\text{B} = \dfrac{x^2 - 4|x| + 4}{|x| - 2} = \dfrac{|x|^2 - 4|x| + 4}{|x| - 2} = \dfrac{(|x| - 2)^2}{|x| - 2} = |x| - 2.\]

\item Η ανίσωση γίνεται:
\begin{align*}
\text{B} - \text{A} < 2d(x, 4) - 5 \ &\Leftrightarrow \\
|x| - 2 - (3 - |x|) < 2|x - 4| - 5 \ &\Leftrightarrow \\
2|x| - 5 < 2|x - 4| - 5 \ &\Leftrightarrow \\
2|x| < 2|x - 4| \ &\Leftrightarrow \\
|x| < |x - 4| \ &\Leftrightarrow \\
|x|^2 < |x - 4|^2 \ &\Leftrightarrow \\
x^2 < x^2 - 8x + 16 \ &\Leftrightarrow \\
8x < 16 \ &\Leftrightarrow \\
x &< 2.
\end{align*}

Δεδομένου ότι για να έχει νόημα η ανίσωση πρέπει $x \ne \pm 1$ και $x \ne \pm 2$, τελικά η ανίσωση
αληθεύει για
\[x \in (-\infty, -2) \cup (-2, -1) \cup (-1, 1) \cup (1, 2).\]
\end{alist}

